% 2011 年广东卷文科第 9 题：由三视图求几何体体积。
% 来源 PDF：普通高考/2011/2011广东文.pdf，SHA-256 b42cd0aaa193fefb6dcb39d1ae186d536a83fafbd4a64975deafbea5613d98ce。
% 原页：第 1 页（A3 三栏）第 2 栏上部，书页码 665。
% 重绘依据：照原页——左上为正（主）视图，等边三角形，左右两顶点分别引出尺寸界线，
%   沿左边外侧以双向箭头标注 2√3；三角形内有一条从顶点到底边中点的竖直线。
%   右上为侧（左）视图，等腰三角形，底边中线向两侧各画一段带箭头的尺寸线，数字 2 落在两箭头之间；
%   三角形内同样有一条竖直线。左下为俯视图，菱形，两条对角线都画出，左上斜边外侧以双向箭头标注 2。
%   三视图下各标名称；原页不画顶点圆点、不标其余尺寸。本图按原页拓扑与标注重绘，不按比例。
\begin{tikzpicture}[scale=0.82, >=stealth, semithick]
  \coordinate (P) at (1.7321,3);
  \coordinate (L) at (0,0);
  \coordinate (R) at (3.4641,0);
  % 正（主）视图：正三角形，顶点到底边中点的竖直线
  \draw (L) -- (R) -- (P) -- cycle;
  \draw (P) -- (1.7321,0);
  % 左边外侧的尺寸界线与尺寸线，标注 2√3
  \draw (L) -- (-0.433,0.25);
  \draw (P) -- (1.299,3.25);
  \draw[<->] (-0.433,0.25) -- (1.299,3.25);
  \node[rotate=60, fill=white, inner sep=1pt] at (0.15,1.98) {$2\sqrt3$};
  \node at (1.7321,-0.72) {正视图};
  % 侧（左）视图：底边 2、高 3 的等腰三角形，含竖直线
  \draw (5.2,0) -- (7.2,0) -- (6.2,3) -- cycle;
  \draw (6.2,3) -- (6.2,0);
  \draw[<->] (5.2,-0.45) -- (5.89,-0.45);
  \draw[<->] (6.51,-0.45) -- (7.2,-0.45);
  \node at (6.2,-0.45) {$2$};
  \draw (5.2,0) -- (5.2,-0.45);
  \draw (7.2,0) -- (7.2,-0.45);
  \node at (6.2,-1.12) {侧视图};
  % 俯视图：对角线 2√3（水平）与 2（竖直）的菱形，两条对角线都画出
  \draw (0,-2.2) -- (1.7321,-1.2) -- (3.4641,-2.2) -- (1.7321,-3.2) -- cycle;
  \draw (0,-2.2) -- (3.4641,-2.2);
  \draw (1.7321,-1.2) -- (1.7321,-3.2);
  % 左上斜边外侧的尺寸界线与尺寸线，标注 2（尺寸线平行于该边、向外侧偏移）
  \draw (0,-2.2) -- (-0.25,-1.767);
  \draw (1.7321,-1.2) -- (1.482,-0.767);
  \draw[<->] (-0.25,-1.767) -- (1.482,-0.767);
  \node[rotate=30, fill=white, inner sep=1pt] at (0.30,-1.55) {$2$};
  \node at (1.7321,-3.88) {俯视图};
\end{tikzpicture}
